\documentclass[11pt]{article}
\usepackage[utf8]{inputenc}
\usepackage[margin=1in]{geometry}
\usepackage{xcolor}
\usepackage{listings}
\usepackage{hyperref}
\usepackage{titlesec}
\usepackage{array}
\usepackage{enumitem}
\usepackage{amssymb}
% Configuration for code snippets (shell/bash style)
\lstset{
    backgroundcolor=\color{gray!10},
    basicstyle=\ttfamily\small,
    breaklines=true,
    frame=single,
    rulecolor=\color{black!30},
    keywordstyle=\color{blue},
    commentstyle=\color{green!50!black}
}
% Hyperlink styling
\hypersetup{
    colorlinks=true,
    linkcolor=blue,
    urlcolor=blue,
    citecolor=blue
}
% Title Formatting
\titleformat{\section}{\normalfont\Large\bfseries}{}{0pt}{}
\titleformat{\subsection}{\normalfont\large\bfseries}{}{0pt}{}
% NOTE TO COURSE STAFF: Replace the Ed Discussion thread name in the
% Submission section with the actual thread you want students posting to
% before publishing.
\begin{document}
% --- HEADER SECTION ---
\begin{center}
    {\Huge \textbf{CS193}} \\
    \vspace{2mm}
    {\Large \textbf{Homework - Week 3}} \\
    \vspace{2mm}
    {\large Exploring Real Weather Data from the Command Line} \\
\end{center}

% --- INTRODUCTION ---
\section{Introduction}
In lecture you learned how to go beyond typing one command at a time and start \textbf{combining commands} in bash --- finding files, and chaining tools together with \textbf{pipes and filters}. This week you'll put those skills to work on a real dataset.
\\\\
The data comes from \textbf{WHIN} (the \href{https://www.whin.org/weather}{Wabash Heartland Innovation Network}), which runs a dense network of weather stations across a 10-county region of north-central Indiana. You'll work with a fixed snapshot of readings from \textbf{24 stations} and answer questions about it entirely from the command line.
\\\\
You only need \textbf{bash and a terminal} --- no other languages, no installs. A helpful companion reference is \textit{The Linux Command Line} (\href{https://linuxcommand.org/tlcl.php}{linuxcommand.org/tlcl.php}), especially the chapters on navigation and I/O redirection.

% --- GETTING STARTED ---
\section{Getting Started}
Download and unzip the \texttt{cs193\_week3} folder. It already contains everything you need:
\begin{itemize}
    \item \texttt{whin\_stations.csv} --- the snapshot, with \textbf{one row per station}. Columns are: \texttt{station\_id}, \texttt{station\_name}, \texttt{county}, \texttt{latitude}, \texttt{longitude}, \texttt{temp\_f}, \texttt{humidity\_pct}, \texttt{wind\_mph}, \texttt{rain\_in}, \texttt{soil\_moist\_pct}, \texttt{observed\_at}.
    \item \texttt{station\_logs/} --- one log file per station, each with that station's hourly readings.
\end{itemize}

% --- HOW TO RUN ---
\section{How to Work Through It}
\begin{enumerate}
    \item Open a terminal (in VS Code: \textbf{Terminal \(\rightarrow\) New Terminal}, or any terminal app).
    \item \texttt{cd} into the project folder so the file names below work:
\begin{lstlisting}[language=bash]
$ cd cs193_week3
\end{lstlisting}
    \item Run each command below, read the output, and record your answer. No server, compiler, or installation is required --- it's plain bash.
\end{enumerate}

% --- PART A ---
\section{Part A --- Finding Your Way Around (~7 min)}
Every station writes its own log file. Let's locate them.

\begin{enumerate}[label=\textbf{A\arabic*.}]
    \item List the log files, then count them. How many stations logged data today?
\begin{lstlisting}[language=bash]
$ ls station_logs
$ ls station_logs | wc -l
\end{lstlisting}

    \item Now count \textit{every} CSV file in the folder, including the snapshot. \texttt{find} searches inside sub-folders for you.
\begin{lstlisting}[language=bash]
$ find . -name "*.csv" | wc -l
\end{lstlisting}

    \item WHIN's Tippecanoe County stations all start with \texttt{TIP}. List just those log files.
\begin{lstlisting}[language=bash]
$ find station_logs -name "TIP*"
\end{lstlisting}

    \item Print the full path of every \textbf{White County} log file (their IDs start with \texttt{WHI}).
\begin{lstlisting}[language=bash]
$ find station_logs -name "WHI*"
\end{lstlisting}
\end{enumerate}

% --- PART B ---
\section{Part B --- Reading the Data with Pipes \& Filters (~10 min)}
Now work with the snapshot file \texttt{whin\_stations.csv}.

\begin{enumerate}[label=\textbf{B\arabic*.}]
    \item Peek at the first 5 lines. What does the very first line contain?
\begin{lstlisting}[language=bash]
$ head -5 whin_stations.csv
\end{lstlisting}

    \item How many stations are in the snapshot? The header row is not a station, so skip it with \texttt{tail -n +2} (``start at line 2'').
\begin{lstlisting}[language=bash]
$ tail -n +2 whin_stations.csv | wc -l
\end{lstlisting}

    \item Print just the station names (column 2). The file is comma-separated, so tell \texttt{cut} the delimiter is a comma.
\begin{lstlisting}[language=bash]
$ cut -d, -f2 whin_stations.csv
\end{lstlisting}

    \item How many stations are in \textbf{Tippecanoe} County? \texttt{grep -c} counts matching lines.
\begin{lstlisting}[language=bash]
$ grep -c Tippecanoe whin_stations.csv
\end{lstlisting}

    \item Count how many stations are in \textit{each} county. Chain four tools: take the county column, sort it, then collapse duplicates with a count.
\begin{lstlisting}[language=bash]
$ tail -n +2 whin_stations.csv | cut -d, -f3 | sort | uniq -c
\end{lstlisting}

    \item Which station is the \textbf{coolest} in this snapshot? Temperature is column 6. Sort the rows numerically by that column and take the first one.
\begin{lstlisting}[language=bash]
$ tail -n +2 whin_stations.csv | sort -t, -k6 -n | head -1
\end{lstlisting}
\end{enumerate}

\textbf{Challenge (optional):} Write a \textbf{single command line} that prints only the station names in \textbf{White} County. Hint: \texttt{grep} to pick the rows, then pipe into \texttt{cut} to keep the name column.

% --- SUBMITTING ---
\section{How to Submit}
\begin{enumerate}
    \item Paste each command you ran and its output into a plain-text file named \texttt{week3\_lastname.txt}.
    \item At the bottom, add one sentence answering: \textit{``Which county has the most stations, and how did you find out?''}
    \item Go to Ed Discussion and open the \textbf{Week 3} submission thread.
    \item Create a new post, attach your \texttt{week3\_lastname.txt}, and submit.
\end{enumerate}

% --- GRADING ---
\section{Grading}
\begin{itemize}
    \item \textbf{A1} --- count the log files with \texttt{ls | wc -l} --- 10 pts
    \item \textbf{A2} --- count all CSVs with \texttt{find ... | wc -l} --- 10 pts
    \item \textbf{A3} --- list \texttt{TIP*} files with \texttt{find} --- 10 pts
    \item \textbf{A4} --- list \texttt{WHI*} files with \texttt{find} --- 10 pts
    \item \textbf{B1} --- \texttt{head}; identifies the first line as the header row --- 8 pts
    \item \textbf{B2} --- skips the header with \texttt{tail -n +2} and counts with \texttt{wc -l} --- 8 pts
    \item \textbf{B3} --- extracts the name column with \texttt{cut -d, -f2} --- 8 pts
    \item \textbf{B4} --- counts Tippecanoe rows with \texttt{grep -c} --- 8 pts
    \item \textbf{B5} --- \texttt{cut | sort | uniq -c} chain for per-county counts --- 10 pts
    \item \textbf{B6} --- \texttt{sort -t, -k6 -n | head -1} for the coolest station --- 8 pts
    \item \textbf{Submission \& reflection} --- correctly named file with commands, outputs, and the reflection sentence --- 10 pts
    \item \textbf{Challenge (bonus)} --- \texttt{grep White | cut -d, -f2} in one line --- +5 pts
\end{itemize}
\textbf{Total: 100 points (+5 bonus).}

\end{document}
